\chap{持续开发与快速交付：在并行变化中推进演化}
\label{sec:delivery}
\subsection{持续变化是工作情境}
多个 Agent 同时工作时，共享代码和依赖可能一直在变化。若每次发布都等待所有工作完成，等待时间会随并行活动增多而上升。智能开发系统需要支持局部成果在具备交付条件时推进，并在之后继续吸收新的变化。持续集成强调频繁集成与反馈\cite{fowlerci}；本文进一步关注由 Agent 自主选择实际交付范围和完成修复的组织方式。

“局部”应依据行为与依赖定义。一个功能可能跨越前端、服务和配置，送达其中一部分会造成不完整行为。相反，同一文件中的多项相关改动有时难以低成本拆开，可以作为一个小批次共同检查与交付。责任成员需要判断本次实际内容能否形成可解释的行为结果。

在任务取消示例中，前端显示新的等待状态依赖后端提供相应语义。如果只发布前端，而后端仍把接收请求解释为完成，界面就无法可靠呈现结果。最小可用批次因此包含状态协议及其消费方；执行器内部与这项协议无关的优化可以留到后续交付。

\subsection{确定性工具与智能判断的分工}
范围判断、依赖理解和异常解释适合由 Agent 结合语义完成；实际内容捕获、并行合成、版本切换和结果记录应由可重复工具处理。把确定性操作交给稳定工具，有助于减少每个成员重新编排低层步骤的成本。工具应返回实际执行结果，使 Agent 能够继续判断，而不是只返回“已开始”。

一次交付可以包含准备、捕获、核对、合成、切换和观察。准备阶段允许不同任务并行，捕获阶段确定本次实际内容，最终合成与切换只在必要的短时间内协调。如果捕获之后又发生修改，后续修改进入下一批交付，不必无限等待代码停止变化。

并行合成需要保留他人已经交付的结果。对互不相关的变化可以直接组合，对同一契约或内容的冲突则应由相关责任成员解释。工具能识别文本冲突，却未必识别语义冲突；两个补丁都能应用，仍可能对状态含义作出互斥假设。定向检查需要覆盖这些实际共享关系。

\fig{06-delivery}{并行开发与小批量交付。各模块持续推进，发布工具捕获一组满足依赖的实际内容；后续变化进入下一批工作，交付后观察与修复继续进行。}{delivery}

\subsection{验证与发布范围相互说明}
交付记录应说明送达了哪些内容、依赖哪些前提，以及进行了什么检查。验证结果也应能定位到实际送达的内容和环境。若测试是在另一份代码上完成，团队需要判断差异是否相关。这个判断可以局部进行，而不必因任何无关变化就宣布全部证据失效。

定向验证的“充分”来自所要支持的结论。例如，一次修改只改变取消等待提示，可能需要检查显示和状态映射；若修改执行器停止过程，则需要检查资源释放与并发完成分支。全量回归可以作为独立活动发现组合问题，触发条件由实际风险和成本决定。发布次数本身不能说明验证是否足够。

旧版本出现故障时，修复交付往往具有更高的时间价值。系统可以通过相同的局部捕获与核对路径送达修复，再观察是否恢复。修复本身也可能失败，因此需要保留实际结果和可恢复入口。快速交付的价值应体现在问题被发现到有效修复的时间，而非动作频率。

\subsection{一次交付的最小事实记录}
一次交付记录可以包含目的、实际内容标识、相关依赖、执行环境、验证入口和最终结果。它回答的是“什么变化实际到达了哪里”，而不只是“原计划发布什么”。如果工具只保存选定路径，却没有保存捕获时的内容，后续代码变化会使排障依据变得含糊。

例如，状态服务与界面共同交付之后，另一成员又修改了同一个界面文件。原记录应仍能找回先前送达的内容，同时把新修改识别为下一批候选。发生故障时，Agent 可以比较实际运行内容与证据，判断问题来自本批变化、共同依赖或后续修改。这样的记录支持局部核对，减少全局重建状态的成本。

交付记录也可以驱动面向人的进展视图。它提供已送达内容与证据入口，Agent 再根据目标说明用户可以体验什么。若切换失败，界面应显示失败及当前运行状态；若只完成准备，则不能把候选版本当成可用成果。工程记录与交互表达共享事实，有助于避免两个层面逐渐产生差异。

\subsection{环境隔离与恢复的条件性取舍}
开发入口和稳定入口可以共享部分数据或运行基础，也可以进行更强隔离。选择取决于系统架构、任务持续时间、数据约束和运维成本。共享有助于减少同步与迁移工作，却会使某些故障同时影响多个入口；隔离有助于限制影响，也会增加状态一致性和部署维护成本。

Portal 的实践讨论倾向于轻量滚动更新与快速修复，并接受部分运行基础共享。这项经验适用于理解设计取舍，不能直接推广为所有项目的发布规则。涉及不可逆数据操作或严格外部承诺的系统，需要相应恢复策略和更强检查。框架提供的共同原则是让约束明确、范围可判断、结果可观察。

代码恢复与数据恢复也需要分别考虑。重新启用旧代码不一定撤销新版本已经写入的数据。协议和数据格式的兼容条件应在变化前被识别，长时间运行的任务则需要明确在版本变化中怎样继续完成。保持任务身份和结果连续性，是判断交付是否成功的一部分。

\subsection{低成本记录与交付后的学习}
保留实际内容、交付目的、检查范围和运行结果，能够帮助后续排障与维护。这类记录可以由工具自动生成大部分内容，再由 Agent 补充必要解释。它们应服务于找回事实与恢复工作，避免把繁重的整理流程放到每次局部交付之前。

交付后的观察可能改变原先判断。一个修改通过代码测试，却使用户更难理解状态；另一个修改功能正确，但频繁引起跨模块同步。团队可以把这些结果反馈到验证策略、模块职责和界面设计中。持续交付因而也是获取工程知识的过程。

\begin{insight}{交付完成的含义}
一次交付至少要能够说明实际送达的内容、已核对的行为和仍未解决的问题。对于重要用户任务，还应提供可体验的结果。代码已经修改、工具已经启动或成员已经汇报完成，都只是其中的部分事实。
\end{insight}

局部交付与长期架构质量应共同观察。若每次快速修复都增加特殊分支，后续变化成本可能持续上升。团队可以定期或在出现明显信号时整理重复问题，判断是否需要调整协议和模块职责。何时进行这类改进，同样可以由 Agent 结合实际收益选择。
